% !TEX root =../main.tex

% PREÁMBULO (Mover a preamble.tex si se desea)
%%%%%%%%%%%%%%%%%%%%%%%%%%%%%%%%%%%%%%%%%%%%%%%%%%%%%%%%%%%%
\documentclass[12pt,a4paper]{article}

\usepackage[]{babel}
\usepackage[utf8]{inputenc}
\usepackage[T1]{fontenc}

\usepackage{graphicx}
\usepackage{amsmath}
\usepackage{caption}
\usepackage{float}
\usepackage{csquotes}
\usepackage{bookmark}
\usepackage{hyperref}
\usepackage{xurl}
\usepackage{listings}
\usepackage{xcolor}
\usepackage{minted}
\usepackage{placeins}
\usepackage{subcaption}
\usepackage{enumitem}
\usepackage{tikz}
\usetikzlibrary{patterns}
\usetikzlibrary{shapes.geometric, arrows.meta, positioning}
\usetikzlibrary{calc}
\usepackage{xstring}


\usepackage{geometry}
\geometry{margin=2.5cm}

\usepackage{hyperref}

% Formato IEEE para referencias
\usepackage[style=ieee, backend=biber]{biblatex}


\addbibresource{./referencias.bib}

\newcommand{\hex}[1]{\texttt{#1}}
\newcommand{\code}[1]{\texttt{#1}}

\definecolor{red}{HTML}{B04A4A}
\definecolor{blue}{HTML}{4A6FA5}
\definecolor{yellow}{HTML}{C9A227}
\definecolor{purple}{HTML}{7A5C99}

% Cisco terminal env
\usepackage[dvipsnames]{xcolor}
\definecolor{strings}{HTML}{ee8888}
\definecolor{comments}{HTML}{ffffff}
\definecolor{prompt}{HTML}{8888ee}
\definecolor{keywords}{HTML}{fcff95}
\definecolor{background}{HTML}{000000}
\definecolor{numbers}{HTML}{a884e0}
\definecolor{rfc}{HTML}{4e85cd}
\definecolor{defaulttext}{HTML}{eeeeee}
 
\usepackage{listings}
\makeatletter
\DeclareRobustCommand{\filename}[1]{
    \begingroup
        \def\textendash{-}
        \filename@parse{#1}
        \edef\filename@base{\detokenize\expandafter{\filename@base}}
        \filename@base.\filename@ext
    \endgroup
}
\makeatother
\lstdefinestyle{default}{
    backgroundcolor=\color{background},
    breakatwhitespace=true,
    breaklines=true,
    basicstyle=\ttfamily\small\color{defaulttext},
    commentstyle=\color{comments},
    keywordstyle=\color{keywords},
    numberstyle=\color{numbers},
    rulecolor=\color{numbers},
    deletekeywords={},
    escapeinside={}{},
    extendedchars=true,
    frame=lines,
    keepspaces=true,
    morekeywords={},
    numbers=left,
    showspaces=false,
    showstringspaces=false,
    showtabs=false,
    stepnumber=1,
    stringstyle=\color{strings},
    tabsize=2,
    title=\filename{\lstname}
}
 
\lstset{
    style=default,
}
 
\lstdefinelanguage{cisco-terminal}{
    basicstyle=\ttfamily\footnotesize\color{prompt},
    morecomment=[l][\color{defaulttext}]{\#},
    morecomment=[l][\color{defaulttext}]{>},
    morecomment=[s][\color{strings}]{'}{'},
    numbers=left
}

\usepackage{tcolorbox}
\tcbuselibrary{listings, skins}
\definecolor{boxbg}{HTML}{000000} % black background
\definecolor{boxborder}{HTML}{FFFFFF} % optional border colo

\newtcblisting{darkcode}[1][]{
    listing only,
    colback=boxbg,       % box background color
    boxrule=0pt,
    title=\texttt{#1},
    coltitle=white,
    colbacktitle=black,
    colframe=boxbg,
    toptitle=3mm,        % padding above title text
    bottomtitle=3mm,     % padding below title text
    listing options={
        language=cisco-terminal,
        basicstyle=\ttfamily\footnotesize\color{prompt},
        commentstyle=\color{comments},
        keywordstyle=\color{keywords},
        numberstyle=\color{numbers},
        rulecolor=\color{numbers},
        deletekeywords={},
        escapeinside={}{},
        extendedchars=true,
        numbers=none,
        breaklines=true,
        keepspaces=true,
        showspaces=false,
        showstringspaces=false,
        showtabs=false,
        tabsize=2,
        title=#1,
        frame=none,
        morekeywords={config}
    },
    colframe=boxborder,  % border color
    left=10pt,            % padding inside the box
    right=10pt,
    top=10pt,
    bottom=10pt,
    arc=10pt,             % slightly rounded corners
}

% swa -> switch A
% swaA -> Switch A en la parte A
\newcommand{\swaA}{\textcolor{blue}{\small\texttt{001d.718a.4700}}}
\newcommand{\swbA}{\textcolor{yellow}{\small\texttt{0026.0af1.6b00}}}
\newcommand{\swcA}{\textcolor{purple}{\small\texttt{001d.7166.f180}}}
\newcommand{\swdA}{\textcolor{red}{\small\texttt{001c.f6b4.2480}}}
% swaB -> Switch A en la parte B
\newcommand{\swaB}{\textcolor{blue}{\small\texttt{046c.9dbe.f500}}}
\newcommand{\swbB}{\textcolor{yellow}{\small\texttt{dcce.c109.5580}}}
\newcommand{\swcB}{\textcolor{purple}{\small\texttt{b4a8.b95d.3180}}}
\newcommand{\swdB}{\textcolor{red}{\small\texttt{001d.7166.f180}}}
